\documentclass{article}

\usepackage{amsmath,amssymb}
\usepackage{xurl}
\usepackage[hidelinks]{hyperref}
\setlength{\emergencystretch}{2em}

% Shared commands for notes in docs/paper-gaps/.

\DeclareUnicodeCharacter{2080}{\ifmmode{}_{0}\else\textsubscript{0}\fi}
\DeclareUnicodeCharacter{2081}{\ifmmode{}_{1}\else\textsubscript{1}\fi}
\DeclareUnicodeCharacter{2082}{\ifmmode{}_{2}\else\textsubscript{2}\fi}
\DeclareUnicodeCharacter{2099}{\ifmmode{}_{n}\else\textsubscript{n}\fi}

\newcommand{\Lean}{\textsc{Lean}}
\ExplSyntaxOn
\NewDocumentCommand{\leanid}{m}{
  \ifmmode
    \textsf{\detokenize{#1}}
  \else
    \group_begin:
    \urlstyle{sf}
    \tl_set:Nn \l_tmpa_tl {#1}
    \tl_replace_all:Nnn \l_tmpa_tl {\_} {_}
    \exp_args:NV \path \l_tmpa_tl
    \group_end:
  \fi
}
\ExplSyntaxOff
\providecommand{\tr}{\operatorname{tr}}
\newcommand{\tnleanIssueBase}{https://github.com/LionSR/TNLean/issues/}
\newcommand{\tnleanPrBase}{https://github.com/LionSR/TNLean/pull/}
\newcommand{\tnleanDocBase}{https://sirui-lu.com/TNLean/docs/}
\newcommand{\ghissue}[1]{\href{\tnleanIssueBase#1}{GitHub issue \##1}}
\newcommand{\ghpr}[1]{\href{\tnleanPrBase#1}{PR \##1}}
\newcommand{\doclink}[1]{\href{\tnleanDocBase#1.html}{\path{#1}}}

% Dirac notation and the identity operator, as in blueprint/src/macros/common.tex.
% \providecommand keeps note-local definitions authoritative.
\usepackage{dsfont}
\providecommand{\ket}[1]{|#1\rangle}
\providecommand{\bra}[1]{\langle #1|}
\providecommand{\braket}[2]{\langle #1 | #2 \rangle}
\providecommand{\ketbra}[2]{|#1\rangle\!\langle #2|}
\providecommand{\Id}{\mathds{1}}
\providecommand{\one}{\mathds{1}}
\providecommand{\C}{\mathbb{C}}
\providecommand{\MN}[1]{M_{#1}(\C)}
\providecommand{\MD}{M_D(\C)}

% External referencing for paper sources.  The build helper writes sanitized
% auxiliary files under build/paper-aux/ and excludes duplicated source labels.
\usepackage{xr}

% Import a generated paper auxiliary file with a caller-supplied label prefix.
% The candidate paths support builds from docs/paper-gaps/, the repository root,
% and build/paper-aux/ itself.
\newcommand{\importpaperaux}[2]{%
  \IfFileExists{../../build/paper-aux/#2.aux}{%
    \externaldocument[#1]{../../build/paper-aux/#2}%
  }{%
    \IfFileExists{build/paper-aux/#2.aux}{%
      \externaldocument[#1]{build/paper-aux/#2}%
    }{%
      \IfFileExists{#2.aux}{%
        \externaldocument[#1]{#2}%
      }{}%
    }%
  }%
}

% General external reference macros
\newcommand{\paperref}[2]{\ref{#1-#2}}
\newcommand{\papereqref}[2]{(\ref{#1-#2})}

% CPSV16 source (1606.00608 / MPDO-22-12-17-2.tex) external document and shortcuts
\importpaperaux{cpsv16-}{1606.00608/MPDO-22-12-17-2-tnlean}
\newcommand{\cpsvref}[1]{\paperref{cpsv16}{#1}}
\newcommand{\cpsveqref}[1]{\papereqref{cpsv16}{#1}}

% Verdict marker: \gapnote{<kind>}{<status>}. Kind is the policy
% classification (clarification, local-correction, scope-restriction,
% unfaithful, false-source, open-gap); status is open, wip (elimination
% underway), resolved, or historical. Machine-read by the site index;
% prints a verdict line.
\newcommand{\gapnote}[2]{\par\noindent\textbf{Verdict:} #1 (#2).\par}


\title{The Physical Complement in the Topological Gibbs Decomposition}
\author{}
\date{2026-07-23}

\begin{document}
\maketitle
\gapnote{open-gap}{resolved}

\section{Status}

The physical-coordinate construction is resolved on the source-defined domain
of chain lengths $L\geq2$. The generic fusion-clause theorem is
\leanid{MPOTensor.BNTFusionTensorClause.hasPhysicalTopologicalGibbsDecomposition}.
The literal CPSV fixed-point capstone is
\leanid{MPOTensor.physicalTopologicalGibbsDecomposition_of_isRFPViaTS_of_cpsvCanonicalForm}.
The earlier normalized BNT-refined corollary
\leanid{MPOTensor.physicalTopologicalGibbsDecomposition_of_isRFPViaTS} remains
available.

These theorems construct a finite extension to the physical complement,
transport the retained projectors, and prove the physical density formula. The
separate undefined length-one convention is recorded in
\path{docs/paper-gaps/cpsv16_topological_gibbs_length_one.tex}.

\section{Source statement}

The theorem following Theorem~4.14 in
Ref.~\cite[lines~1013--1016]{Cirac2016MPDO_arXiv} states the physical-chain
identity
\begin{align}
    \rho^{(N)}(M)
        &= \sum_{i=1}^{d}\lambda_iP_i^{(N)}e^{-H_N},
    \label{eq:cpsv16_topological_gibbs_physical_complement_source_decomposition}\\
    [P_i^{(N)},e^{-H_N}]
        &= 0.
    \label{eq:cpsv16_topological_gibbs_physical_complement_source_commutation}
\end{align}
The operators $P_i^{(N)}$ are orthogonal projections on the physical Hilbert
space, and
\begin{align}
    H_N
        &= \sum_j h_{j,j+1}
    \label{eq:cpsv16_topological_gibbs_physical_complement_source_hamiltonian}
\end{align}
is a translationally invariant commuting nearest-neighbor Hamiltonian on that
space.

The derivation in Ref.~\cite{Cirac2016MPDO_arXiv} first applies the vertical
canonical form from Proposition~4.13 and then suppresses the resulting local
change of coordinates. Formally, this change of coordinates is a retained-row
coisometry
\begin{align}
    V
        &\colon \mathbb C^d\longrightarrow\mathbb C^s,
    \label{eq:cpsv16_topological_gibbs_physical_complement_vertical_map}\\
    VV^\dagger
        &= I_s,
    \label{eq:cpsv16_topological_gibbs_physical_complement_coisometry}\\
    V^\dagger V
        &= E.
    \label{eq:cpsv16_topological_gibbs_physical_complement_retained_projection}
\end{align}
The projection $E$ need not equal $I_d$, because the original physical space
may contain a zero sector. This point is documented in
\path{docs/paper-gaps/cpgsv17_vertical_isometry_zero_sector.tex}.

\section{The retained-coordinate theorem}

Let $\mu$ be the strictly positive diagonal multiplicity matrix on
$\mathbb C^s$, and define the diagonal logarithmic multiplicity energy by
\begin{align}
    k
        &= -\log\mu.
    \label{eq:cpsv16_topological_gibbs_physical_complement_retained_energy}
\end{align}
For every $N\geq2$, the retained-space interaction is
\begin{align}
    h
        &= k\otimes I_s.
    \label{eq:cpsv16_topological_gibbs_physical_complement_retained_interaction}
\end{align}
Its periodic translates commute. Since each site occurs once as the first
site of a translated term,
\begin{align}
    \exp\left(-\sum_j h_{j,j+1}\right)
        &= \mu^{\otimes N}.
    \label{eq:cpsv16_topological_gibbs_physical_complement_retained_exponential}
\end{align}

Let $R_N$ denote the retained-chain operator, let
$P_{i,\mathrm{ret}}^{(N)}$ be its terminal spectral projections, and let
$H_{N,\mathrm{ret}}$ be the retained Hamiltonian. The retained decomposition
is
\begin{align}
    R_N
        &= \sum_i\lambda_iP_{i,\mathrm{ret}}^{(N)}
           e^{-H_{N,\mathrm{ret}}}.
    \label{eq:cpsv16_topological_gibbs_physical_complement_retained_decomposition}
\end{align}
The terminal projections commute with $e^{-H_{N,\mathrm{ret}}}$. Applying
$(V^\dagger)^{\otimes N}$ on the left and $V^{\otimes N}$ on the right of
Eq.~\ref{eq:cpsv16_topological_gibbs_physical_complement_retained_decomposition}
reconstructs $\rho^{(N)}(M)$ exactly.

This reconstruction does not yet prove the physical-chain statement of
Ref.~\cite[lines~1013--1016]{Cirac2016MPDO_arXiv}. The Hamiltonian
$H_{N,\mathrm{ret}}$ and the projections
$P_{i,\mathrm{ret}}^{(N)}$ act on $(\mathbb C^s)^{\otimes N}$, whereas the
source theorem requires operators on $(\mathbb C^d)^{\otimes N}$.

\section{Failure of direct compression}

Direct transport of the one-site retained Gibbs factor gives
\begin{align}
    G_{\mathrm{comp}}
        &= V^\dagger e^{-k}V.
    \label{eq:cpsv16_topological_gibbs_physical_complement_direct_compression}
\end{align}
The matrix $G_{\mathrm{comp}}$ vanishes on $(I_d-E)\mathbb C^d$. If
$E\neq I_d$, it is singular and therefore cannot equal the exponential of a
finite matrix on the full physical one-site space.

A finite extension begins with
\begin{align}
    k_{\mathrm{phys}}
        &= V^\dagger kV,
    \label{eq:cpsv16_topological_gibbs_physical_complement_physical_energy}\\
    e^{-k_{\mathrm{phys}}}
        &= V^\dagger e^{-k}V+(I_d-E).
    \label{eq:cpsv16_topological_gibbs_physical_complement_physical_gibbs_factor}
\end{align}
The complementary term assigns finite Gibbs weight one to
$(I_d-E)\mathbb C^d$.

For $L\geq2$, let $H_{L,\mathrm{phys}}$ be the sum of the periodic translates
of $k_{\mathrm{phys}}\otimes I_d$. Its Gibbs factor is
\begin{align}
    e^{-H_{L,\mathrm{phys}}}
        &= \bigotimes_{j=1}^{L}
           \left(V^\dagger e^{-k}V+I_d-E\right).
    \label{eq:cpsv16_topological_gibbs_physical_complement_physical_chain_extension}
\end{align}
This identity differs from the false global equality
\begin{align}
    e^{-H_{L,\mathrm{phys}}}
        &=
        (V^\dagger)^{\otimes L}
        e^{-H_{L,\mathrm{ret}}}
        V^{\otimes L}.
    \label{eq:cpsv16_topological_gibbs_physical_complement_false_global_transport}
\end{align}
When $E\neq I_d$, the right-hand side of
Eq.~\ref{eq:cpsv16_topological_gibbs_physical_complement_false_global_transport}
vanishes on every complementary or mixed sector, whereas the left-hand side
has finite positive weight there. Thus the displayed equality is false in
that case.

\section{Transported physical projectors}

Define the physical projectors by
\begin{align}
    \widehat P_i^{(N)}
        &= (V^\dagger)^{\otimes N}
           P_{i,\mathrm{ret}}^{(N)}
           V^{\otimes N}.
    \label{eq:cpsv16_topological_gibbs_physical_complement_transported_projectors}
\end{align}
The coisometry relation
$VV^\dagger=I_s$ implies their projection and orthogonality identities. Each
$\widehat P_i^{(N)}$ is supported on $E^{\otimes N}$ and therefore annihilates
all complementary and mixed-sector contributions in
Eq.~\ref{eq:cpsv16_topological_gibbs_physical_complement_physical_chain_extension}.
Consequently, the retained commutation relations imply
\begin{align}
    [\widehat P_i^{(N)},e^{-H_{N,\mathrm{phys}}}]
        &= 0.
    \label{eq:cpsv16_topological_gibbs_physical_complement_transported_commutation}
\end{align}
Exact vertical reconstruction and the retained decomposition then give
\begin{align}
    \rho^{(N)}(M)
        &= \sum_i\lambda_i\widehat P_i^{(N)}
           e^{-H_{N,\mathrm{phys}}}.
    \label{eq:cpsv16_topological_gibbs_physical_complement_physical_decomposition}
\end{align}
Thus
Eqs.~\ref{eq:cpsv16_topological_gibbs_physical_complement_transported_commutation}
and
\ref{eq:cpsv16_topological_gibbs_physical_complement_physical_decomposition}
supply the physical Hamiltonian and physical projections required by the
source statement.

\section{Resolution}

The formalization resolves the physical-complement gap through four steps:
\begin{enumerate}
    \item It proves
    Eq.~\ref{eq:cpsv16_topological_gibbs_physical_complement_physical_gibbs_factor}
    from $VV^\dagger=I_s$.
    \item It defines the fixed physical two-site interaction
    $k_{\mathrm{phys}}\otimes I_d$ and proves
    Eq.~\ref{eq:cpsv16_topological_gibbs_physical_complement_physical_chain_extension}
    for every chain length at least two.
    \item It transports the retained terminal projections by
    Eq.~\ref{eq:cpsv16_topological_gibbs_physical_complement_transported_projectors}
    and proves their projection, orthogonality, and commutation identities.
    \item It combines these identities with exact vertical reconstruction to
    obtain
    Eq.~\ref{eq:cpsv16_topological_gibbs_physical_complement_physical_decomposition}.
\end{enumerate}

The generic theorem
\leanid{MPOTensor.BNTFusionTensorClause.hasPhysicalTopologicalGibbsDecomposition},
the CPSV fixed-point theorem
\leanid{MPOTensor.physicalTopologicalGibbsDecomposition_of_isRFPViaTS_of_cpsvCanonicalForm},
and the normalized BNT-refined corollary
\leanid{MPOTensor.physicalTopologicalGibbsDecomposition_of_isRFPViaTS} formalize
this construction.

The result applies on the source-defined domain $N\geq2$. The separate
one-site ambiguity in the periodic two-site interaction remains recorded in
\path{docs/paper-gaps/cpsv16_topological_gibbs_length_one.tex}; the
physical-complement construction neither resolves nor supplies that convention.

\bibliographystyle{alpha}
\bibliography{references}

\end{document}
